
\tikzset{%
  every neuron/.style={
    circle,
    draw,
    %fill=black,
    radius=0.1mm
  },
  neuron missing/.style={
    draw=none,
    scale=2,
    text height=0.25cm,
    execute at begin node=\color{black}$\vdots$
  },
}

\begin{center}

\begin{tikzpicture}%[ultra thick]




  \node [every  neuron ] (u-1) at ( 10 ,-0) { } ;
  \node [every  neuron ] (u-2) at ( 8 ,-1.5) { } ;
  \node [every  neuron ] (u-3) at ( 12 ,-1.5) { } ;
  \node [every  neuron ] (u-4) at ( 6.5 ,-3) { } ;
  \node [every  neuron ] (u-5) at ( 9.5 ,-3) { } ;
  \node [every  neuron ] (u-6) at ( 10.5 ,-3) { } ;
  \node [every  neuron ] (u-7) at ( 13.5 ,-3) { } ;
%  \node [every  neuron ] (u-8) at ( 0.5 ,-4.5) { } ;
%  \node [every  neuron ] (u-9) at ( 1.5 ,-4.5) { } ;
%  \node [every  neuron ] (u-10) at ( 2.5 ,-4.5) { } ;
%  \node [every  neuron ] (u-11) at ( 1.25 ,-6) { } ;
%  \node [every  neuron ] (u-12) at ( 1.75 ,-6) { } ;
%  \node [every  neuron ] (u-3) at ( 6 ,-3) { } ;
%\node (u-4) at ( 0.0 ,-6) { $ \vdots $  };
%
%\node [every  neuron ] (u-6) at ( 5 ,-6) { $l$  };
%\node [every  neuron ] (u-7) at ( 7 ,-6) { $r$  };
%\node  (u-8) at ( 5 ,-7) { $\vdots$  };
%\node  (u-9) at ( 7 ,-7) { $\vdots$  };
\draw [->] (u-1) -- (u-2);
\draw [->] (u-1) -- (u-3);
\draw [->] (u-2) -- (u-4);
\draw [->] (u-2) -- (u-5);
\draw [->] (u-3) -- (u-6);
\draw [->] (u-3) -- (u-7);
%\draw [->] (u-4) -- (u-8);
%\draw [->] (u-5) -- (u-9);
%\draw [->] (u-5) -- (u-10);
%\draw [->] (u-9) -- (u-11);
%\draw [->] (u-9) -- (u-12);





  \node [every  neuron ] (v-1) at ( 5 ,-0) { } ;
  \node [every  neuron ] (v-2) at ( 4 ,-1.5) { } ;
  \node [every  neuron ] (v-3) at ( 3 ,-3) { } ;
  \node [every  neuron ] (v-4) at ( 2 ,-4.5) { } ;
  \node [every  neuron ] (v-5) at ( 1 ,-6) { } ;
  \node [every  neuron ] (v-6) at ( 0 ,-7.5) { } ;
%  \node [every  neuron ] (v-7) at ( 6 ,-3) { } ;
%  \node [every  neuron ] (v-8) at ( 0.5 ,-4.5) { } ;
%  \node [every  neuron ] (v-9) at ( 1.5 ,-4.5) { } ;
%  \node [every  neuron ] (v-10) at ( 2.5 ,-4.5) { } ;
%  \node [every  neuron ] (v-11) at ( 1.25 ,-6) { } ;
%  \node [every  neuron ] (v-12) at ( 1.75 ,-6) { } ;
%  \node [every  neuron ] (v-3) at ( 6 ,-3) { } ;
%\node (v-4) at ( 0.0 ,-6) { $ \vdots $  };
%
%\node [every  neuron ] (v-6) at ( 5 ,-6) { $l$  };
%\node [every  neuron ] (v-7) at ( 7 ,-6) { $r$  };
%\node  (v-8) at ( 5 ,-7) { $\vdots$  };
%\node  (v-9) at ( 7 ,-7) { $\vdots$  };
\draw [->] (v-1) -- (v-2);
\draw [->] (v-2) -- (v-3);
\draw [->] (v-3) -- (v-4);
\draw [->] (v-4) -- (v-5);
\draw [->] (v-5) -- (v-6);
%\draw [->] (v-3) -- (v-7);
%\draw [->] (v-4) -- (v-8);
%\draw [->] (v-5) -- (v-9);
%\draw [->] (v-5) -- (v-10);
%\draw [->] (v-9) -- (v-11);
%\draw [->] (v-9) -- (v-12);
%%  node[right=6pt] {$y$} at
%\draw[decoration={brace,raise=6pt},decorate, thick]
%(8,-1.5) -- node[right=10pt, align=left] {  $n^{\prime}$-size binary tree, \\It's guaranteed that any node, \\beside the root, \\satisfies the heap inequality. } (8, -7.5);
%
%\draw[draw=gray, dashed] (4.5,-7.5) rectangle ++(3.5,6);
\end{tikzpicture} 


\end{center}
